\documentclass[11pt]{article}
\usepackage[T1]{fontenc}
\usepackage{lmodern,amsmath,amssymb,microtype}
\usepackage[margin=1.02in]{geometry}
\usepackage[colorlinks=true,urlcolor=blue,citecolor=blue]{hyperref}
\setlength{\emergencystretch}{2em}
\hypersetup{pdfauthor={Wentao Li},pdftitle={Deciding efficient generation for pairs of deterministic context-free languages}}
\title{Deciding efficient generation for pairs of deterministic context-free languages}
\author{Wentao Li\\[2pt]\small Prepared with assistance from OpenAI Codex}
\date{November 2026}
\begin{document}
\maketitle
\begin{abstract}
For a pair of deterministic context-free languages with infinite intersection, an unrestricted computable generator can always make zero mistakes. Deciding whether a polynomial-time generator can do this is $\Sigma^0_2$-complete. The same classification holds if one allows randomization and an arbitrary finite uniform bound on the expected number of mistakes. A fixed pair can force linearly many expected mistakes on arbitrarily long compatible prefixes. These statements concern explicit string output.
\end{abstract}

\paragraph{Setting.} Two languages over a fixed finite alphabet are given by deterministic pushdown automata. A target presents its strings in any complete enumeration without repetition. After seeing distinct strings $x_1,\ldots,x_j$ the generator writes a string. It makes a mistake if it repeats a presented string or misses the target. The time per round is bounded by a fixed polynomial in $S_j=1+\sum_{i\le j}(1+|x_i|)$. A polynomial and an algorithm may depend on the given pair. We allow internal memory and, for randomized algorithms, independent fair bits. A uniform mistake bound is one finite bound across both targets, all texts and all horizons. If the intersection is infinite, exhaustive search in it always finds an unseen word.

\paragraph{Theorem.} Even under the promise of infinite intersection, each of the following decision problems on a pair of deterministic pushdown automata is $\Sigma^0_2$-complete (computable many-one reductions with promised outputs).
\begin{enumerate}
\item Does a deterministic polynomial-time zero-mistake generator exist?
\item Does a randomized polynomial-time generator with a finite uniform expected mistake bound exist?
\end{enumerate}
Here $\Sigma^0_2$ consists of predicates expressible as $\exists u\,\forall v\,R(u,v)$ with decidable $R$. For membership on the promise, guess a clocked algorithm, polynomial time bound, and, where needed, an integer mistake bound. Universally test finite distinct prefixes and each possible target. Membership in the target is decidable and every valid finite prefix extends to a complete enumeration. The algorithm's finite computation tree gives the exact expected number of mistakes up to that prefix. Uniform boundedness over horizons is equivalent to the bound for complete texts by monotone convergence. This gives a $\Sigma^0_2$ predicate agreeing with each decision problem on its promise.

\paragraph{Two automata for a clock.} Given an index $e$, simulate one transition per phase of an effective enumerator of the c.e. set $W_e$. Arrange that it emits a pulse on completing each previously unseen enumerated element, by keeping a list and performing the duplicate-removal computation within the simulated enumerator. The pulses are finite in number exactly when $W_e$ is finite. Build a deterministic Turing machine with a unary positive integer $q$, initially two. Have a unique checkpoint state, visited at the start and after every phase, with an outgoing transition. Following the next simulation transition, replace $q$ by $q+1$ for an ordinary phase, or $2^q$ on a pulse. Sweep the complete new unary block before the checkpoint. Repeated doubling $q$ times constructs $2^q$ in time polynomial in the new value. Ordinary work and individual simulation steps can likewise be implemented in polynomial time in the new value. A fixed finite-tape construction suffices, since the tape alphabet and number of tapes can be absorbed in finite configuration symbols. Alternatively a polynomial-overhead deterministic one-tape simulation with a distinguished state reached once per simulated checkpoint gives a one-tape clock, which we use below.

Let $w_j$ be the entire computation history through checkpoint $j$, encoded as follows. Configurations use one marked state/head position and tape symbols, with the rightmost padding extending by one cell after each step (a one-sided tape suffices). Give the initial row a fixed width containing the initial tape. Alternate the direction of encoding of successive rows and include row separators and an end marker. There is then a unique encoded trace to each checkpoint. The language of well-formed candidate histories with the designated initial configuration and a checkpoint in the last row is regular: require one state/head per row, track row parity, and remember the type of the last state until the separator or end marker.

One deterministic pushdown automaton additionally verifies transitions with even starting index, and a second verifies those with odd starting index, numbering the initial row zero. For a checked pair it pushes the first row and reads the next row against the popped symbols, aligned by the reversal convention. A single transition changes only the scanned cell, state and neighbouring head position. Thus the aligned symbol pairs for a correct step form a regular language, in either reading direction: outside a window of at most two cells the symbols agree, and the window must realize a transition from a finite table. Check the extra rightmost cell against a virtual old blank, first or last according to the parity, and require simultaneous row end after this offset. This handles a head moving onto the added cell. For unchecked boundary rows only check syntax, and permit a final pushed row to end unpaired. The separators, row parity and end marker determine each action without guessing. Denote the languages by $L_0,L_1$. Their intersection is exactly $\{w_1,w_2,\ldots\}$ since all successive transitions and both boundary conditions are checked. If desired give language $i$ a private regular branch starting with its own fresh letter. This preserves determinism and the intersection. The construction is effective, and all finite alphabets can use a fixed-length binary code uniformly within each instance.

\paragraph{The size gap.} Write $q_j$ for the value at checkpoint $j$. Since $j\le q_j$, monotonicity and the phase time bound imply that the cumulative running time at checkpoint $j$ is bounded by a fixed polynomial in $q_j$. The time-padded trace has quadratic size in this time up to constant factors. On the common prefix $w_1,\ldots,w_j$ the full input size $S_j$ is therefore polynomial in $q_j$. On a pulse the next trace has at least $2^{q_j}$ symbols from the sweep alone. All later traces are strictly longer. These bounds tolerate the single-tape and binary encodings.

If $W_e$ is finite then $q_j=j+O_e(1)$ eventually. On an arbitrary history of $j$ examples, compute the first $j+1$ checkpoint words and output one not yet presented. This costs a polynomial in $S_j$, including all finite exceptions by increasing the constant. It makes zero mistakes.

If $W_e$ is infinite, fix any polynomial-time generator. On all sufficiently late pulse rounds of the common prefix it cannot write a fresh word in the intersection within its clock. A symbol-by-symbol output model is being used here. On each such round the sum of the two mistake indicators on the identical prefix is at least one, for every random tape. The two runs can be coupled with identical tapes since their observed examples are identical. Up through any $k$ sufficiently late pulses, some target has at least $k/2$ expected mistakes on a prefix extendible to a complete text for that target. Thus no finite uniform expected bound exists. The classical $\Sigma^0_2$-complete index set $\mathrm{FIN}=\{e:W_e\text{ is finite}\}$ reduces to both questions. For completeness, given $\exists u\,\forall v\,R(e,u,v)$, output at stage $s$ the least $u\le s$ passing all tests $v\le s$, or $s+1$ if none does. The values stabilize if there is a true witness and escape every bound otherwise. They enumerate a finite set exactly in the affirmative case. This proves completeness.\hfill$\square$

Using a pulse in every phase also gives a single fixed pair on which the expected lower bound through round $T$ is $(T-O_G(1))/2$ on one of the two extensions. In contrast, for a fixed finite family of regular languages with infinite common intersection there is a polynomial-time zero-mistake algorithm. Its fixed product automaton has an accepting path with a nonempty pumpable loop. Pumping it $j+1$ different numbers of times gives $j+1$ distinct common words of length $O(j+1)$, one of which is fresh.

\paragraph{Scope and priority.} X. Li, A. Han, D. Shi, J. Jiang and J. Gao pose hardness for natural language classes and structural criteria for efficient closure generation in §5.2, \url{https://arxiv.org/abs/2608.05029}. Our theorem gives an unconditional standard running-time obstruction for deterministic context-free languages. Its cause is the length of the next safe output. Their oracle-query model permits arbitrary local computation, so our theorem does not transfer to that cost measure. M. Arenas, P. Barcel\'o, L. Cofr\'e and A. Kozachinskiy use split computation-history checks with a finite intersection for their sample-complexity lower bound, \url{https://proceedings.mlr.press/v306/arenas26a.html}. Here the intersection is always infinite and the distribution of checkpoint lengths encodes finiteness. Independent proof and priority review are pending.
\end{document}
